\scriptsize
\begin{longtable}{@{}p{2.2cm}p{1.0cm}p{4.7cm}p{7.0cm}@{}}
\caption{The sixteen rating questions as asked, with their five levels (0 to 4, separated by slashes). Dimension and form are those of the second study. \textsuperscript{n}: asked in the second study only.}\label{tab:wording}\\
\toprule
Question & Dim., form & As asked & Levels 0 to 4 \\
\midrule
\endfirsthead
\toprule
Question & Dim., form & As asked & Levels 0 to 4 \\
\midrule
\endhead
\bottomrule
\endfoot
Persuasion & F1, A & How strongly does the speaker try to persuade the audience to accept a viewpoint or take an action? & No persuasive appeal. / A slight persuasive cue. / A moderate persuasive appeal. / Strong or repeated persuasive appeals. / Persuasion is sustained and central throughout. \\
\addlinespace[2pt]
Persistence & F1, B & How strongly does the speaker persist in explaining or supporting a point? & No persistence cue. / A slight continuation cue. / Moderate effort to continue or support a point. / Strong, repeated effort to develop a point. / Very sustained effort to develop or reinforce a point. \\
\addlinespace[2pt]
Positive tone & F2, A & Where does the overall emotional tone conveyed by the speaker's own language fall from negative to positive? Rate expressed tone, not underlying feeling. & Clearly negative overall tone. / Somewhat negative overall tone. / Neutral or balanced/mixed tone. / Somewhat positive overall tone. / Clearly positive overall tone. \\
\addlinespace[2pt]
Disappointment & F2, B & How strongly does the speaker's own language convey disappointment? & No disappointment cue. / A slight cue. / A moderate cue. / Strong or repeated disappointment. / Very prominent, sustained disappointment. \\
\addlinespace[2pt]
Self-disclosure & F3, B & How much personal experience or personal information does the speaker disclose about themselves? & No personal disclosure. / A brief, minor disclosure. / A moderate personal disclosure. / Substantial or repeated personal disclosure. / Extensive personal disclosure throughout. \\
\addlinespace[2pt]
Self-critical & F3, A & How strongly does the speaker's own wording criticize or devalue themselves? & No self-critical cue. / A slight self-critical cue. / A moderate self-critical statement. / Strong or repeated self-criticism. / Very prominent, sustained self-criticism. \\
\addlinespace[2pt]
Compassion & F4, B & How strongly does the speaker's own wording convey compassion or concern for others? & No compassion cue. / A slight cue of compassion. / A moderate cue of compassion. / Strong or repeated compassionate wording. / Very prominent, sustained compassion. \\
\addlinespace[2pt]
Respect & F4, A\textsuperscript{n} & How strongly does the speaker's own wording convey respect toward people discussed? & No respect cue. / A slight cue of respect. / A moderate cue of respect. / Strong or repeated respectful wording. / Very prominent, sustained respect. \\
\addlinespace[2pt]
Narrative & F5, B & How strongly is the excerpt organized as a narrative with events or a sequence over time? & No narrative sequence. / A slight narrative cue. / A partial sequence of events. / A clear, developed event sequence. / A strongly developed narrative structure throughout. \\
\addlinespace[2pt]
Rebuttal & F5, A\textsuperscript{n} & How strongly does the speaker answer or rebut an opposing claim? & No rebuttal. / A slight response to opposition. / A moderate rebuttal. / Strong or repeated rebuttals. / Rebuttal is sustained and central throughout. \\
\addlinespace[2pt]
Energetic & F6, A & How strongly does the wording itself convey verbal energy or animation? Do not use audio volume or pitch. & No verbal-energy cue. / Slightly animated wording. / Moderately animated wording. / Strongly animated or energetic wording. / Exceptionally energetic wording throughout. \\
\addlinespace[2pt]
Calm & F6, B & How strongly does the wording convey a calm, even-handed presentation? & No calmness cue. / A slight calmness cue. / A moderate calm and measured presentation. / Strongly calm and measured wording. / Exceptionally steady and measured wording throughout. \\
\addlinespace[2pt]
Humor & F7, A & How strongly does the speaker use humor or joking in their own narration? & No humor cue. / A slight joke or humorous cue. / A moderate use of humor. / Strong or repeated humor. / Humor is exceptionally prominent throughout. \\
\addlinespace[2pt]
Imaginative & F7, B\textsuperscript{n} & How strongly does the speaker use imaginative, novel, or creative framing? & No imaginative framing. / A slight imaginative cue. / A moderate creative or novel framing. / Strong or repeated imaginative framing. / Very prominent, sustained creative framing. \\
\addlinespace[2pt]
Hedging & F8, B & How strongly does the speaker qualify, hedge, or signal uncertainty about their own statements? & No qualification cue. / A slight qualification cue. / A moderate amount of qualification. / Frequent or strong qualification. / Very pervasive qualification or uncertainty marking. \\
\addlinespace[2pt]
Knowledge limits & F8, A\textsuperscript{n} & How strongly does the speaker distinguish what they know from what they do not know or cannot confirm? & No explicit knowledge-boundary cue. / A slight qualification. / A moderate distinction or caveat. / Clear and repeated knowledge limits. / Knowledge boundaries are explicit and consistently maintained. \\
\end{longtable}
